\section{Using impulses}
We can improve the performance of the previous feedback cancellation
solution if we take also into consideration that the \textbf{feedback}
signal $\tilde{s}$ that finally reaches our microphone is (at least in
part) the
\href{https://en.wikipedia.org/wiki/Convolution}{convolution} of
$s$ and a signal $h$ that represents the impulse response of our
local audioset (speaker, mic, walls, monitor, keyboard, our body, ...)
to an \href{https://en.wikipedia.org/wiki/Impulse_response}{impulse
  signal $\delta(t)$}.\footnote{This technique is similar to the
  carried out by submarines when they user the sonar to perform
  echo-location, or by bats to fly in the darkness.} In other words,
we can modify Eq.~\eqref{eq:simplest} to compute
\begin{equation}
 \tilde{\mathbf{n}}[t] = \mathbf{m}[t] - (\mathbf{h}\ast\mathbf{s}[t-d]) = \mathbf{m}[t] - (\mathbf{h}\ast\mathbf{s})[t-d],
  \label{eq:using_convolution}
\end{equation}
where $\ast$ represents the (digital) convolution between (in our case
of) digital signals, and ${\mathbf h}$ is the digitalized version of
$h(t)$.

Notice that in this solution, Eq.~\eqref{eq:minimal_filter} (the
estimation for the echo signal) becomes
\begin{equation}
  \hat{\mathbf s}[t] \triangleq (\mathbf{h}\ast\mathbf{s})[t-d].
  \label{eq:echo_signal_modeled_by_FIR}
\end{equation}

\subsection{Filtering in the frequency domain}
The convolution of digital signals in the time domain can be expensive
(with
\href{https://en.wikipedia.org/wiki/Computational_complexity_theory}{computational
  complexity} $O^2$, where $O$ is the number of elements to process)
if the number of samples or/and filter coefficientsis is
high. Fortunately, thanks to the
\href{https://en.wikipedia.org/wiki/Convolution}{convolution theorem}
\cite{kovacevic2013fourier,Oppenheim2}, the convolution can be
replaced by the \href{https://en.wikipedia.org/wiki/Dot_product}{dot
  product} (with complexity $O$), when we consider the signals in the
frequency domain. Thanks to this, we can rewrite the
Eq.~\eqref{eq:using_convolution} as
\begin{equation}
  \tilde{\mathbf n}[t] = {\mathbf m}[t] - ({\mathcal F}^{-1}\{{\mathbf S}{\mathbf H}\})[t-d],
  \label{eq:faster}
\end{equation}
where ${\mathbf S}$ is the (digital) Fourier transform\footnote{The
Fourier transform is an special case of the Laplace transform where
$\sigma=0$ in the complex (Laplace) domain represented by
$s=\sigma+j\omega$ frequencies. This simplification can be used for
the characterization of our near-end audioset because it can be
considered as a FIR (Finite Impulse Response) system (in ausence of an
audio signal, the echo signal always decays with the time). Finally,
notice that ${\mathbf S}[\omega]$ (the $\omega$-th frequency component
of ${\mathbf S}$) is a complex number (the Fourier coefficients are
complex numbers).} of ${\mathbf s}$, ${\mathbf H}$ is the Fourier
transform\footnote{For computing the Fourier transform of digital
signals use a library such as
\href{https://numpy.org/doc/2.1/reference/routines.fft.html}{\text{numpy.fft}}.}
of ${\mathbf h}$, and ${\mathcal F}^{-1}$ represents the inverse
(digital) Fourier transform. Notice that all these transforms are
applied to digital signals, and there exist fast algorithms (with
complexity $O\log_2O$) to ``travel'' from the signal domain to the
frecuency domain, and viceversa.

\subsubsection*{Practical issues}
Unfortunately, even when we expect that this improved echo
supression algorithm is going to perform better than the previous one,
the computation of the filter weights ${\mathbf h}$ requires emitting
impulses that can be heard by the user. Furthermore, while the
acoustic response of the environment is being analyzed, the near-end
should remain silent; otherwise, its own voice would be treated as an
echo and attenuated by the filter.

%\section{Adaptive filtering using impulses}
%Unfortunately, none of the previous solutions is ``dynamic'', in the
%sense that if we modify the audioset, some parameters must be modified
%in real-time:
%\begin{enumerate}
%\item In Equations \eqref{eq:simplest} and
%  \eqref{eq:using_convolution}, the parameter $d$ which expresses (in
%  sample-times) the time that $s$ need to travel from the speaker to
%  the microphone, depends, for example, on the orientation of the
%  screen of the laptop.
%\item Eq.~\eqref{eq:using_convolution} and (obviously)
%  Eq~\eqref{eq:faster} are also affected by $h$, which depends on
%  current audioset configuration.
%\end{enumerate}

\begin{comment}
Filters can be implemented in the frequency domain (such as happens in
Eq.~\eqref{eq:faster}) or in the signal (time in our case)
domain. Signal domain (digital) convolutions are affordable if the
length\footnote{The number of coefficients.} of the (digital) filters
is small.\footnote{Take in mind that convolution is a $O^2$ operation,
  and therefore, we can only handle in real-time with our computers
  small filter.}

Eq.~\eqref{eq:minimal_filter} represents the minimal expression for a
\href{https://en.wikipedia.org/wiki/Finite_impulse_response}{FIR
  filter} (whose spectral shape is completely flat because it only has
a single coefficient and therefore, it attenuates equally all the
frequencies). In general, the expression of a FIR filter in the
time-domain (applied to a signal that previously has been delayed $d$
samples) is
\begin{equation}
  \hat{\mathbf s}[t] = \sum_{k=0}^{K-1}{\mathbf h}_k{\mathbf s}[t-k-d],
  \label{eq:FIR_filter}
\end{equation}
where $K$ is the number of coefficients of the filter ${\mathbf h}$,
and ${\mathbf h}_k$ is its $k$-th coefficient. When the (coefficients
of the) filter are updated for each input sample,
Eq.~\eqref{eq:FIR_filter} can be rewritten as
\begin{equation}
  \hat{\mathbf s}[t] = \sum_{k=0}^{K-1}{\mathbf h}_k^{(t)}{\mathbf s}[t-k-d].
  \label{eq:adaptive_FIR_filter}
\end{equation}

In general, the filter should not be updated so often (e.g., every
second could be fast enough). In this case
Eq.~\eqref{eq:adaptive_FIR_filter} should be
\begin{equation}
  \hat{\mathbf s}[t] = \sum_{k=0}^{K-1}{\mathbf h}_k^{(i)}{\mathbf s}[t-k-d].
  \label{eq:adaptive_FIR_filter_second}
\end{equation}
where $i$ represents the update iteration.
If we decide to work with impulses, ${\mathbf h}^{(i)}$ is determined
by the samples recorded by our soundcard after $d$ sample-times of the
generation of each impulse when the impulses should be the only source
of sound. If there were other sources of sound (such as for example,
our voice) the updates should be discarded or averaged using some
\href{https://en.wikipedia.org/wiki/Moving_average}{moving-average}
technique.

Finally, notice that we must send
\begin{equation}
  \tilde{\mathbf n}[t] = \mathbf{m}[t] - \hat{\mathbf s}[t].
  \label{eq:finally}
\end{equation}
\end{comment}

\begin{comment}
\section{Subtraction in the frequency domain}
Ultimately, feedback in InterCom occurs because the audio signal we
are generating is reproduced through our speakers (with a certain time
lag). In this feedback situation, the spectrum of a newly received
chunk is similar to the spectrum of a newly sent chunk. Therefore, if
the chunk we want to send is pre-filtered (before being sent) using
the opposite filter to that described by the spectrum of the played
chunk, we should cancel the feedback.

Let ${\mathbf M} = {\mathcal F}({\mathbf m})$ the Fourier transform of
the $i$-th chunk captured by the soundcard, and let
${\mathbf S} = {\mathcal F}({\mathbf s})$ the Fourier transform of the
$i$-th played chunk. Considering the reasoning followed in the
previous paragraph, and considering also the Eq.~\eqref{eq:simplest}, it can be concluded that
\begin{equation}
  \tilde{\mathbf N} = {\mathbf M} - a{\mathbf S},
  \label{eq:simplest_fourier}
\end{equation}
where $\hat{\mathbf{N}} = \mathcal{F}(\hat{\mathbf{n}})$, the Fourier
transform of the chunk that should be sent in the $i$-th iteration
(chunk).
\end{comment}

\begin{comment}
Notice that if you prefer to estimate ${\mathbf h}$ using a ``train''
(a sequence) of impulses, the temporal distance between them would 

  For estimating the feedback signal at the
  sample-time $t$, ${\mathbf f}[t]$, the length of a
  \href{https://en.wikipedia.org/wiki/Finite_impulse_response}{FIR
    filter} should be at least $d$, because at least we need $d$
  samples to detect the feedback signal. Therefore, we have that
\begin{equation}
  \hat{\mathbf s}[t] = \sum_{k=0}^{d-1}{\mathbf h}_k^{(t)}{\mathbf s}[t-k],
  \label{eq:LMS_feedback}
\end{equation}
where ${\mathbf h}$ is the near-end impulse response in the time
domain.
\end{comment}

\begin{comment}
\section{Filtering in the frequency domain}

Assume that $s$ (the signal received by our microphone from our
speaker) has essentially the same frequency spectrum as the audio
signal being played (i.e., $h(t)=\delta(t)$, and therefore, the
environment does not modify the spectrum of the reproduced signal). We
denote the signal reproduced by our speaker as $f(t)$ ($f$ for
far-end). In this case, we can propose a feedback reduction system by
applying a filter $\mathbf{F}^{-1}$ to the captured signal
$\mathbf{s}$, which is equal to the ``inverse'' of the spectrum of
$\mathbf{f}$ (the signal we are reproducing). Specifically, this would
involve:

\begin{enumerate}
\item \textbf{Calculating} $\mathbf{F}$, the Fourier transform of the
  reproduced audio signal $\mathbf{f}$, chunk by chunk.
\item \textbf{Calculating} $\mathbf{F}^{-1}$, the ``inverse'' of
  $\mathbf{F}$. If $w_{\text{max}}$ is the frequency component of
  $\mathbf{F}$ with the highest energy, $g_{\text{max}}$, i.e.,
  $\mathbf{F}[w_{\text{max}}]=g_{\text{max}}$, and
  $\mathbf{F}[w_{\text{min}}]=g_{\text{min}}$, normalize its spectrum
  with:
  \begin{equation}
    \mathbf{F}_{\text{normalized}} = \frac{\mathbf{F} - g_{\text{min}}}{g_{\text{max}} - g_{\text{min}}},
  \end{equation}
  and finally calculate:
  \begin{equation}
    \mathbf{F}^{-1} = 1 - \mathbf{F}_{\text{normalized}}.
  \end{equation}
\item \textbf{Filtering the signal captured by the microphone} using
  $\mathbf{F}^{-1}$, taking into account the attenuation factor $a$,
  which relates how much of the reproduced signal is captured by the
  microphone (how much of $\mathbf{f}$ is in $\mathbf{s}$), i.e.,
  calculate:
  \begin{equation}
    \hat{\mathbf{N}} = a\mathbf{F}^{-1}\mathbf{S},
  \end{equation}
  where $\mathcal{F}(\mathbf{s})=\mathbf{N}$.
\item \textbf{Estimating the attenuation factor} $a$ for the next
  chunk. Initially, we assume that $a=1$ (and therefore, that there is
  no attenuation). If it is an overestimation, then some frequency
  components in the sent audio could disappear, and if we
  underestimate, then some frequency components could be
  amplified. Therefore, the attenuation factor should be high enough
  so that the energy of $\mathbf{s}$ does not increase progressively,
  but low enough so that some frequency components do not become 0.
\end{enumerate}
The signal to send (without the feedback) is:
\begin{equation}
  \mathbf{n} = \mathcal{F}^{-1}(\mathbf{N})
\end{equation}

\end{comment}

\begin{comment}
\subsection{Considering the frequency respose of the near-end audioset}
In general, Eq.~\eqref{eq:echo_problem} is oversimplified because our
local audioset (speaker, mic, walls, our body, etc.) does not have a
flat response in the frequency domain. Therefore, a more realistic
model of the echo effect is
\begin{equation}
   {\mathbf m}[i] = {\mathbf n}[i] + \{f({\mathbf s})\}[i-d],
  \label{eq:more_realistic_echo_problem}
\end{equation}
where $f({\mathcal s})$ is a filtered version of the signal
${\mathbf s}$, in which some frequencies are partially attenuated
(notice that now there is a possiblely different $a$-value for each
frequency of ${\mathbf s}$).

Now the problem (apart from finding $d$, obviously) is how to
determine $f(\cdot)$. Again, a simple way of doing this is to use a
call signal to measure the frequency response of the near-end
audioset. For this, a good call signal can be a sequence of
impulses\footnote{A impulse generates a flat
  \href{https://en.wikipedia.org/wiki/Spectral_density}{spectrum}~\cite{kovacevic2013fourier,Oppenheim2}
  for a short period of time.} of a sequence of uniform random noise
signals\footnote{Uniform random noise (or
  \href{https://en.wikipedia.org/wiki/White_noise}{white noise}) has a
  flat spectrum.}. Notice that the frequency response (the filter
coefficients of $f(\cdot)$) can be found comparing the flat spectrum
signal ${\mathbf s}$ generated by the speaker with the supposedly
non-flat spectrum of the signal ${\mathbf m}$ captured by the
microphone. To apply $f(\cdot)$ we can multiply
${\mathcal F}({\mathbf s})$ by the filter coefficients of $f(\cdot)$
in the Fourier domain, being ${\mathcal F}(\cdot)$ the Fourier
Transform~\cite{kovacevic2013fourier,Oppenheim2}.


Lets $H(s)$ the representation of $h(t)$ in the frequency (Laplace)
domain, where $s=\sigma+j\omega$ is a complex number that denotes a
frequency component.

If we consider that $h(t)$ in the frequency domain is denoted by
$H(s)$, (when $s=\sigma+j\omega$ we say that $H(s)$ is the Laplace
Transform of $h(t)$), that $S(s)$ is the representation in the
frequency domain of the signal $s(t)$, and that in our specific case,
where our near-end audioset is a FIR (Finite Impulse Response)
system\footnote{In ausence of an audio signal, the echo always decays
  with the time.}, we can also characterize it using the Fourier
Transform (where $sigma=0$), i.e., we can say that $H(j\omega)$ (the
Fourier Transform of $h(t)$) describes our local audioset from a
frequency perspective. Considering now the convolution theorem in the
Fourier domain, we can rewrite the Eq.~\eqref{eq:simplest} as
\begin{equation}
  \tilde{\mathbf n}[i] = {\mathbf m}[i] - {\mathbf h}({\mathbf S}{\mathbf H})[i-d],
  \label{eq:simplest2}
\end{equation}
where ${\mathbf S}$ is the Fast Discrete Fourier Transform of
${\mathbf s}$ and ${\mathbf H}$ is the Fast Discrete Fourier Transform
of ${\mathbf h}$, and ${\mathcal F}^{-1}$ represents the Fast Inverse
Discrete Fourier Transform.

Therefore, as it can be seen in Eq.~\eqref{eq:simplest2}, to find the
echo signal we need to mulplity  ${\mathbf S}$ (the played signal in the Fourier
domain) by ${\mathbf H}$ (the frequency response of the near-end
audioset).

\end{comment}
